\documentclass[10pt,a4paper]{amsart}
\usepackage[margin=19mm]{geometry}
\usepackage{amsmath,amssymb,mathtools}
\usepackage[scr=rsfs,cal=euler]{mathalfa}
\usepackage{xfrac}
\usepackage[unicode,hidelinks,bookmarks=false]{hyperref}
\newcommand{\ie}{i.e.\ }
\newcommand{\cf}{cf.\ }
\newcommand{\resp}{respectively\ }
\newcommand{\Subsection}{\subsection}
\providecommand{\nobreakdash}{\nobreak\hbox{-}}
\setlength{\emergencystretch}{2em}
\multlinegap0pt
\usepackage{mathtools}
\usepackage{lmodern}
\usepackage{enumitem}
\numberwithin{equation}{section}
\newtheorem{theorem}{Theorem}[section]
\newtheorem{proposition}[theorem]{Proposition}
\newtheorem{lemma}[theorem]{Lemma}
\newtheorem{corollary}[theorem]{Corollary}
\newtheorem{question}[theorem]{Question}
\newtheorem{remark}[theorem]{Remark}
\newcommand{\C}{\mathbb{C}}
\newcommand{\R}{\mathbb{R}}
\newcommand{\Z}{\mathbb{Z}}
\newcommand{\PP}{\mathbb{P}}
\newcommand{\cO}{\mathcal{O}}
\newcommand{\dbar}{\bar\partial}
\newcommand{\half}{\tfrac12}
\newcommand{\iso}{\xrightarrow{\sim}}
\DeclareMathOperator{\Ric}{Ric}
\DeclareMathOperator{\Aut}{Aut}
\DeclareMathOperator{\Kur}{Kur}
\DeclareMathOperator{\Alb}{Alb}
\DeclareMathOperator{\Spec}{Spec}
\DeclareMathOperator{\Vol}{Vol}
\begin{document}
\title[A negative K"ahler-Einstein threefold with non-integrable in]{A negative K\"ahler-Einstein threefold with non-integrable infinitesimal Einstein deformations}
\author{Ari Krishna}
\date{}
\maketitle
\noindent\emph{Source and licence.} A negative K\"ahler-Einstein threefold with non-integrable infinitesimal Einstein deformations. Version arXiv:2608.13481v1 (CC BY 4.0). This material is distributed under the Creative Commons Attribution 4.0 International licence, http://creativecommons.org/licenses/by/4.0/, as recorded in the arXiv OAI licence metadata for this exact version and shown on the abstract page https://arxiv.org/abs/2608.13481v1. Full rights notice: licenses/arxiv-2608.13481-CC-BY-4.0.txt.\par\smallskip
\noindent\textit{Editorial scope.} Counted unit: the original \texttt{lemma} environment \texttt{lem:descent} at source lines 126--145 together with its complete proof at lines 147--166; the delivered document keeps source lines 88--182 so that every referenced label and every citation resolves inside it. Native editable source is the paper's own concatenated TeX; the AMS wrapper is editorial formatting only. No publisher class, upstream script or unrelated artwork is redistributed.
\section{Constructing the threefold}

We write $q(S)=h^1(S,\cO_S)$.

\begin{proposition}\label{prop:surface}
There exist a smooth projective surface $S$, a faithful action of $G\cong(\Z/7)^4$
on $S$, and a class $\alpha\in H^1(S,T_S)^G$ such that
$$K_S\ \text{is ample},\qquad q(S)=0,\qquad
H^1(S,T_S)=\C\alpha,$$
and
$$
[\alpha,\alpha]\neq0\quad\text{in }H^2(S,T_S).$$
\end{proposition}

\begin{proof}
B\"ohning, Graf von Bothmer, and Pignatelli construct a smooth surface $S$ with
$$\Kur(S)\cong\Spec\C[x]/(x^2),$$
$K_S$ ample, and $q(S)=0$ \cite[Thm.~1.1 and Thms.~6.3, 6.5]{BVP}. Their surface is a totally ramified abelian cover whose deck group is $G\cong(\Z/7)^4$ \cite[Cor.~5.4]{BVP}. The natural transformation of \cite[Lem.~2.11]{BVP} identifies the tangent space of the branch-data deformation functor with $H^1(S,T_S)^G$. The verification in the proof of \cite[Thm.~5.5]{BVP}, using \cite[Thm.~3.19]{BVP}, of the hypotheses of \cite[Cor.~2.13]{BVP} shows that this transformation is an isomorphism. Hence,
$$H^1(S,T_S)=H^1(S,T_S)^G.$$

The aforementioned description of $\Kur(S)$ gives $\dim_{\C}H^1(S,T_S)=1$; let $\alpha$ be a generator. Since $K_S$ is ample, $H^0(S,T_S)=0$, so the deformation functor is prorepresentable. The first-order deformation associated with $\alpha$ is induced by
$$\C[x]/(x^2)\longrightarrow\C[t]/(t^2),
\qquad x\longmapsto at,$$
for some $a\neq0$. A lift to $\C[t]/(t^3)$ would have to send $x$ to $at+bt^2$, but
$$(at+bt^2)^2=a^2t^2\neq0\pmod{t^3}.$$
It follows that $\alpha$ has no second-order lift. By the primary obstruction criterion recalled in Lemma~\ref{lem:primary}, this is equivalent to $[\alpha,\alpha]\neq0$.
\end{proof}

Let $B$ be a smooth projective curve of genus $2$. Its fundamental group surjects onto $G$: in the standard presentation, send the four generators of $\pi_1(G)$ to a basis of the abelian group $G$; the surface-group relation maps to zero. By the Riemann existence theorem, the resulting connected topological cover carries a unique structure of a finite \'etale Galois cover
$$q\colon C\longrightarrow B$$
with Galois group $G$. Then, Riemann-Hurwitz gives that the genus of $C$ is 
$$g(C)=1+|G|\bigl(g(B)-1\bigr)=1+7^4=2402.$$
Note that the diagonal action of $G$ on $S\times C$ is free. With this in hand, let
$$X=(S\times C)/G$$
and denote the quotient map by $\pi\colon S\times C\to X$.

We turn to examining the deformation theory of this quotient. Let $p_S$ and $p_C$ be the two projections.


\begin{lemma}[Descent and splitting]\label{lem:descent}
There are natural isomorphisms
$$H^1(X,T_X)\cong H^1(S\times C,T_{S\times C})^G
\cong H^1(S,T_S)^G\oplus H^1(C,T_C)^G.
$$
Moreover, \'etale descent along $q$ identifies
$$
H^1(C,T_C)^G\cong H^1(B,T_B).
$$
If $\alpha_X\in H^1(X,T_X)$ is defined by
$$
\pi^*\alpha_X=p_S^*\alpha,
$$
then, for $\lambda\in\C$ and $\beta\in H^1(B,T_B)$,
\begin{equation}\label{eq:bracket-splitting}
\pi^*\bigl[\lambda\alpha_X+\beta,\lambda\alpha_X+\beta\bigr]
=\lambda^2p_S^*[\alpha,\alpha].
\end{equation}
In particular, the bracket on the left-hand side is nonzero if and only if $\lambda\neq0$.
\end{lemma}

\begin{proof}
Pullback identifies the Kodaira-Spencer differential graded Lie algebra of $X$ with the invariant sub-DGLA on $S\times C$. Since the group-invariants functor under a finite group is exact in characteristic zero, we get
$$H^i(X,T_X)\cong H^i(S\times C,T_{S\times C})^G.$$
Now,
$$T_{S\times C}=p_S^*T_S\oplus p_C^*T_C.$$
The K\"unneth formula, in conjunction with $H^0(S,T_S)=H^0(C,T_C)=0$, yields
$$H^1(S\times C,T_{S\times C})
=H^1(S,T_S)\oplus H^1(C,T_C).$$
The asserted identification for $C\to B$ arises from the same \'etale descent argument.

Let $\widetilde\beta\in H^1(C,T_C)^G$ be the pullback of $\beta$. Classes pulled back from different factors have zero Kodaira-Spencer bracket. Moreover,
$$[\widetilde\beta,\widetilde\beta]=0,$$
because it is represented by a $(0,2)$-form on a curve. This proves \eqref{eq:bracket-splitting}. Finally, the K\"unneth decomposition
\begin{align*}
H^2(S\times C,T_{S\times C})={}&H^2(S,T_S)
\oplus\bigl(H^1(S,T_S)\otimes H^1(C,\cO_C)\bigr)\\
&\oplus\bigl(H^1(C,T_C)\otimes H^1(S,\cO_S)\bigr)
\end{align*}
shows that $p_S^*\colon H^2(S,T_S)\to H^2(S\times C,T_{S\times C})$ is injective. Since $[\alpha,\alpha]$ is $G$-invariant, it remains nonzero after invariant descent. The last assertion follows.
\end{proof}


\section{From complex obstructions to Einstein obstructions}

We first recall the primary complex obstruction. For a compact complex manifold $Z$, its Kodaira-Spencer DGLA is
$$
L_Z=\bigl(A^{0,\bullet}(Z,T_Z),\dbar,[\ ,\ ]\bigr).$$
It governs the Kuranishi deformation functor via the Maurer-Cartan equation
$$\dbar\phi+\half[\phi,\phi]=0;$$
consult \cite{Kodaira,Kuranishi,ManettiDGLA,Sernesi}for details.

\begin{lemma}[Primary obstruction]\label{lem:primary}
Let $Z$ be a compact complex manifold and $\xi\in H^1(Z,T_Z)$. The class $\xi$ admits a second-order lift if and only if
$$[\xi,\xi]=0\quad\text{in }H^2(Z,T_Z).$$
In particular, a class with nonzero bracket is not integrable. The construction can be shown to be functorial under morphisms of Kodaira-Spencer DGLAs.
\end{lemma}


\begin{thebibliography}{99}
\bibitem[BVP]{BVP} C.~B\"ohning, H.-C.~Graf von Bothmer, and R.~Pignatelli, \emph{A rigid, not infinitesimally rigid surface with $K$ ample}, Boll. Unione Mat. Ital. \textbf{15} (2022), 57-85.
\bibitem[Kodaira]{Kodaira} K.~Kodaira, \emph{Complex Manifolds and Deformation of Complex Structures}, Grundlehren der mathematischen Wissenschaften \textbf{283}, Springer, 1986.
\bibitem[Kuranishi]{Kuranishi} M.~Kuranishi, \emph{New proof for the existence of locally complete families of complex structures}, in: Proc. Conf. Complex Analysis (Minneapolis, 1964), Springer, 1965, pp.~142-154.
\bibitem[ManettiDGLA]{ManettiDGLA} M.~Manetti, \emph{Lectures on deformations of complex manifolds}, Rend. Mat. Appl. (7) \textbf{24} (2004), 1-183.
\bibitem[Sernesi]{Sernesi} E.~Sernesi, \emph{Deformations of Algebraic Schemes}, Grundlehren der mathematischen Wissenschaften \textbf{334}, Springer, 2006.
\end{thebibliography}
\end{document}
